\section{Infix to Prefix Conversion}
\label{sec:sq-infix-to-prefix}

\problemstatement{Given an infix expression, convert it to prefix
notation.}

\begin{patternbox}
\emph{``Convert to prefix''} immediately after having just solved
``convert to postfix'' (Section~\ref{sec:sq-infix-to-postfix}) is an
invitation to \textbf{reuse} that algorithm via a symmetry trick, rather
than deriving a new one: prefix is postfix \emph{read backward}, with
operators and operands swapped in scanning direction -- so reversing the
input, adapting the tie-breaking rule, running the same postfix algorithm,
and reversing the output again gets us there with no new derivation.
\end{patternbox}

\subsection*{Understanding the problem carefully}

If postfix writes operators \emph{after} their operands, prefix writes
them \emph{before} -- which is exactly what you get by mirror-imaging a
postfix expression (reversing it character by character) \emph{and}
reversing each multi-character operand if there were any. So the natural
strategy is: reverse the input expression (also swapping every
\texttt{'('} with \texttt{')'} and vice versa, since reversal flips which
parenthesis opens and which closes), run essentially the same
Shunting-Yard scan as Section~\ref{sec:sq-infix-to-postfix} on this
reversed-and-swapped string, then reverse the resulting output once more
to undo the initial reversal.

\begin{ideabox}[The One Rule That Must Change]
When running the postfix algorithm on the reversed string, the tie-break
rule for \textbf{equal precedence} must flip: do \emph{not} pop an
equal-precedence operator already on the stack before pushing a new one,
\emph{regardless of true associativity}. This is because reversing the
scan direction also reverses which operand of a left-associative chain we
encounter first; treating every tie as if it were right-associative during
this reversed scan is what correctly preserves the original
left-to-right grouping once everything is reversed back.
\end{ideabox}

\subsection*{Why reverse-scan-reverse is correct}

This is a direct application of the same ``reverse twice to flip an order
without breaking internal structure'' idea seen in
Section~\ref{sec:sq-queue-using-stack}. Reversing the infix string
converts a left-to-right scan that would naturally produce postfix into
one that, after the tie-break adjustment, produces the mirror image of
postfix -- which, when reversed back to normal reading order, is exactly
prefix. The parenthesis-swapping step is necessary because reversing a
string turns every originally-opening parenthesis into a closing one
(textually) and vice versa, and the algorithm needs them to still behave
correctly as openers and closers during the reversed scan.

\subsection*{Algorithm}

\begin{enumerate}
  \item Reverse $s$; swap every \texttt{'('} with \texttt{')'} and vice
        versa.
  \item Run the Shunting-Yard postfix algorithm
        (Algorithm~\ref{sec:sq-infix-to-postfix}) on this modified
        string, with the tie-break rule changed to \emph{never} pop on
        equal precedence.
  \item Reverse the resulting output string. This is the prefix
        expression.
\end{enumerate}

\subsection*{Dry run}

\dryrun{Take $s = \texttt{"a+b*c"}$.}

\begin{itemize}
  \item Reverse and swap brackets (no brackets here): $\texttt{"c*b+a"}$.
  \item Run the postfix algorithm (never popping on ties):
  \begin{itemize}
    \item \texttt{c}: output $=\texttt{"c"}$.
    \item \texttt{*}: push. Stack $=[*]$.
    \item \texttt{b}: output $=\texttt{"cb"}$.
    \item \texttt{+}: top is \texttt{*} (higher precedence), pop:
          output $=\texttt{"cb*"}$. Push \texttt{+}. Stack $=[+]$.
    \item \texttt{a}: output $=\texttt{"cb*a"}$.
    \item End: pop \texttt{+}: output $=\texttt{"cb*a+"}$.
  \end{itemize}
  \item Reverse the output: $\texttt{"+a*bc"}$.
\end{itemize}

Result: \texttt{"+a*bc"}, correctly representing $a+(b*c)$ in prefix.

\subsection*{C++ implementation}

\begin{lstlisting}
#include <bits/stdc++.h>
using namespace std;

int precedence(char op) {
    if (op == '^') return 3;
    if (op == '*' || op == '/') return 2;
    if (op == '+' || op == '-') return 1;
    return 0;
}

string infixToPrefix(string s) {
    reverse(s.begin(), s.end());
    for (char& c : s) {
        if (c == '(') c = ')';
        else if (c == ')') c = '(';
    }

    string output;
    stack<char> st;

    for (char c : s) {
        if (isalnum(c)) {
            output += c;
        } else if (c == '(') {
            st.push(c);
        } else if (c == ')') {
            while (st.top() != '(') {
                output += st.top();
                st.pop();
            }
            st.pop();
        } else { // operator: never pop on equal precedence
            while (!st.empty() && st.top() != '(' &&
                   precedence(st.top()) > precedence(c)) {
                output += st.top();
                st.pop();
            }
            st.push(c);
        }
    }
    while (!st.empty()) {
        output += st.top();
        st.pop();
    }

    reverse(output.begin(), output.end());
    return output;
}

int main() {
    cout << infixToPrefix("a+b*c") << endl;
    return 0;
}
\end{lstlisting}

\begin{complexitybox}
\textbf{Time Complexity.} $O(n)$, identical in structure to
Section~\ref{sec:sq-infix-to-postfix}.
\textbf{Space Complexity.} $O(n)$.
\end{complexitybox}

\begin{takeawaybox}
``Reverse, adapt the tie-break rule, run the algorithm you already have,
reverse again'' is a powerful general technique for converting a
left-to-right algorithm into its mirror-image counterpart, without
re-deriving the core logic. Recognize when a new problem is the
\emph{mirror image} of one you have already solved before writing new
code from scratch.
\end{takeawaybox}
